\exercisegroup

\begin{enumerate}
\def\labelenumi{\arabic{enumi})}
\tightlist
\item
  设 \(S\) 是任意一个无穷集合。
\end{enumerate}

\begin{enumerate}
\def\labelenumi{\alph{enumi})}
\tightlist
\item
  证明：赋范空间 \(\mathcal{B}(\mathrm{S})\)（由 S 到 \(\mathbf{R}\) 中的有界映射组成；I, 第 4 页, 例）中任一完全集的最小基数等于 \(2^{\mathrm{Card}(\mathrm{S})}\)（考察 S 的诸子集的特征函数所成的集，并注意 \(\mathbf{R}\) 中存在一个处处稠密的可数集）。b) 证明：赋范空间 \(\ell^1 (\mathrm{S})\)（I, 第 4 页, 例）中任一完全集的最小基数等于 \(\mathrm{Card}(\mathrm{S})\)。
\end{enumerate}

\begin{enumerate}
\def\labelenumi{\arabic{enumi})}
\setcounter{enumi}{1}
\item
  在域 R 上的乘积拓扑向量空间 \(E = R^{N}\) 中，用 \(e_{n} (n \in \mathbf{N})\) 表示直和 \(\mathbf{R}^{(\mathbf{N})}\) 的典范基的诸元素。置 \(a_{0} = e_{0}\)，\(a_{n} = e_{0} + (1/n)e_{n}\)（对 \(n \geqslant 1\)）。证明：对每个整数 \(n \geqslant 0\)，诸 \(a_{i}\)（其中 \(0 \leqslant i \leqslant n\)）在 E 中构成一个拓扑无关族，但无穷族 \((a_{n})_{n \geqslant 0}\) 不是拓扑无关的。若 M 是闭向量子空间 \(Ra_{0}\)，则 \(\dot{a}_{n}\)（\(a_{n}\) 在 E/M 中的类，对 \(n \geqslant 1\)）构成一个拓扑无关族，但由诸 \(a_{n}\)（指标 \(n \geqslant 1\)）在 E 中生成的闭向量子空间 N 包含 M。
\item
  设 E 是 R 上的一个拓扑向量空间，f 是 E 的加法群到 R 中的一个同态。证明：若存在 E 中 0 的一个邻域，使 f 在其中有界，则 f 是 E 中的连续线性型。特别地，当 f 半连续（下半连续或上半连续）时即属此种情形。
\item
  用 K 表示赋有绝对值 \(p(\xi) = |\xi|^{1/2}\) 的域 R。设 E 是 K 上的向量空间，由定义在 I = \([0, 1]\) 上、处处右连续且在点 1 处为零的实值受控函数组成；证明：映射 \(x \mapsto \|x\| = \int_{0}^{1} |x(t)|^{1/2} dt\) 在 E 上是一个范数。证明：对 E 中每个函数 \(x \geqslant 0\)，在 E 中存在两个函数 \(x_{1} \geqslant 0\)、\(x_{2} \geqslant 0\)，使得 \(x = \frac{1}{2}(x_{1} + x_{2})\) 且 \(\|x_{1}\| = \|x_{2}\| = \frac{1}{\sqrt{2}} \|x\|\)。由此推出，E 上每个连续线性型都恒等于零。
\item
  设 K 是一个 Hausdorff 拓扑除环，其拓扑是局部逆有界的 (GT, III, § 6, 习题 22)。把 I, 第 12 页 的 prop. 2 与 I, 第 13 页 的 th. 1 推广到 K 上的拓扑向量空间；当 K 还是完备的时候，同样把 I, 第 13 页 的 th. 2 与 I, 第 14 页 的 prop. 3 加以推广。
\item
  设 K 是把 \(Q^{2}\) 的通常拓扑转移到域 \(\mathbf{Q}(\sqrt{2})\) 上（经由映射 \((x, y) \mapsto x + y\sqrt{2}\)）而得到的拓扑除环。
\end{enumerate}

\begin{enumerate}
\def\labelenumi{\alph{enumi})}
\item
  设 \(\mathbf{E}\) 是集合 \(\mathbf{Q}(\sqrt{2})\)，带其 \(\mathbf{K}\) 上的向量空间结构以及由 \(\mathbf{R}\) 的拓扑诱导的拓扑。证明 \(\mathbf{E}\) 是一个 Hausdorff 拓扑向量空间，在 \(\mathbf{K}\) 上的维数为 1，但它不同构于 \(\mathbf{K}_s\)。
\item
  设 \(\mathbf{F}\) 是拓扑向量空间 \(\mathbf{E} \times \mathbf{E}\)（在 \(\mathbf{K}\) 上）；在 \(\mathbf{F}\) 中，超平面 \(\mathbf{E} \times \{0\}\) 是闭的，但不存在连续线性型 \(f\)（在 \(\mathbf{E} \times \mathbf{E}\) 上），使该超平面由方程 \(f(x) = 0\) 给出。
\end{enumerate}

\begin{enumerate}
\def\labelenumi{\arabic{enumi})}
\setcounter{enumi}{6}
\item
  设 K 是一个非离散且非完备的赋值除环，E 是拓扑向量子空间 \(K + Ka\)（\(\hat{K}\) 中，其中 \(a \notin K\)），F 是乘积空间 \(K \times E\)。在 F 中，子空间 \(M = K \times \{0\}\) 是闭的且余维数为 2。设 N 是 F 中由向量 (0, 1) 与 (1, a) 生成的 M 的补子空间；证明 F 不是 M 与 N 的拓扑直和。
\item
  设 \(p\) 是一个素数，\(\mathbf{Q}_p\) 是 \(p\) 进数域 (GT, III, § 6, 习题 23)。设 \(\mathrm{E}_0\) 是拓扑乘积空间 \(\mathbf{Q}_p \times \mathbf{R}\)；若用 \(\mathbf{K}\) 表示带离散拓扑的域 \(\mathbf{Q}\)，则 \(\mathrm{E}_0\) 是 \(\mathbf{K}\) 上的拓扑向量空间。设 \(\mathbf{M}\) 是由元素 \((r, r)\)（其中 \(r\) 取遍 \(\mathbf{Q}\)）组成的向量子空间；又设 \(\theta\) 是一个无理数，\(\mathbf{N}\) 是由元素 \((0, r\theta)\)（其中 \(r\) 取遍 \(\mathbf{Q}\)）组成的向量子空间。设 \(\mathbf{E}\) 是子空间 \(\mathbf{M} + \mathbf{N}\)（在 \(\mathrm{E}_0\) 中）；证明 \(\mathbf{N}\) 是 \(\mathbf{E}\) 中的一个闭超平面，但不存在 \(\mathbf{N}\) 的拓扑补子空间（注意 \(\mathbf{M}\) 在 \(\mathrm{E}_0\) 中处处稠密）。
\item
  设 \(\mathbf{X}\) 是一个 Hausdorff 拓扑空间，\(\mathbf{V}\) 是一个维数 \(n\) 有限的向量子空间，含于空间 \(\mathcal{C}(\mathbf{X};\mathbf{R})\) 内。
\end{enumerate}

\begin{enumerate}
\def\labelenumi{\alph{enumi})}
\item
  证明：存在 \(n\) 个两两不相交的开集 \(\mathrm{U}_i\)（\(1 \leqslant i \leqslant n\)），位于 \(\mathbf{X}\) 中，使得若有函数 \(f \in \mathbf{V}\) 在每个 \(\mathrm{U}_i\) 内都恒等于零，则它在 \(\mathbf{X}\) 中也恒等于零（利用 A, II, § 7.5, cor. 3）。
\item
  设 \(x_{i} \in \mathbf{U}_{i}\)（\(1 \leqslant i \leqslant n\)）。由 a) 推出：存在一个常数 \(c > 0\)，使得对每个函数 \(f \in \mathbf{V}\)，有
\end{enumerate}

\[
\sup _ {x \in \mathrm{X}} | f (x) | \leqslant c \sum_ {i = 1} ^ {n} | f (x _ {i}) |.
\]

\begin{enumerate}
\def\labelenumi{\arabic{enumi})}
\setcounter{enumi}{9}
\tightlist
\item
  设 \(\mathbf{K}\) 是一个局部紧的非离散赋值除环，\(\mathbf{E}\) 是 \(\mathbf{K}\) 上的一个有限维左向量空间。用 \(\mathfrak{N}(\mathbf{E})\) 表示 \(\mathbf{E}\) 上全体范数所成的集，它是 \(\mathcal{C}(\mathbf{E};\mathbf{R})\)（由 \(\mathbf{E}\) 到 \(\mathbf{R}\) 中的（在典范拓扑下）连续映射组成的空间）的一个子空间。
\end{enumerate}

\begin{enumerate}
\def\labelenumi{\alph{enumi})}
\item
  当给 \(\mathcal{C}(\mathrm{E};\mathbf{R})\) 赋以紧收敛拓扑 \(*\)（对此拓扑它是一个 Fréchet 空间）时，集合 \(\mathfrak{N}(\mathrm{E})\) 在 \(\mathcal{C}(\mathrm{E};\mathbf{R})\) 中是闭的，并且是局部紧的。
\item
  设 \(p_0\) 是 \(\mathfrak{N}(\mathrm{E})\) 的一个元素；证明：存在连续映射 \((\lambda, p) \mapsto \pi_{\lambda}(p)\)（\([0, 1] \times \mathfrak{N}(\mathrm{E})\) 到 \(\mathfrak{N}(\mathrm{E})\) 中的），使得 \(\pi_0(p) = p\) 且 \(\pi_1(p) = p_0\)（对每个 \(p \in \mathfrak{N}(\mathrm{E})\)）。
\end{enumerate}

\begin{enumerate}
\def\labelenumi{\arabic{enumi})}
\setcounter{enumi}{10}
\tightlist
\item
  在 I, 第 23 页, 习题 7 的假设下，证明：若 K 与 E 都完备，则 E 的每个闭子空间都有拓扑补（按同处 a) 的方式处理）。
\end{enumerate}

¶ 12) 设 K 是一个局部紧的赋值除环，其绝对值是非离散且超度量的。我们把 K 上左向量空间 E 上满足超度量不等式 (II, 第 2 页) 的范数称为超范数。

\begin{enumerate}
\def\labelenumi{\alph{enumi})}
\tightlist
\item
  设 E 是 K 上的一个有限维左向量空间，\(\alpha\) 是 E 上的一个超范数，H 是 E 中由方程 \(\langle x, a^{*} \rangle = 0\) 给出的一个超平面。证明：存在一点 \(x_{0} \in E\)，函数 \(x \mapsto |\langle x, a^{*} \rangle| / \alpha(x)\) 在 \(E \setminus \{0\}\) 上于该点达到其上确界；证明此时
\end{enumerate}

\[
\alpha (x) = \sup \left(\alpha \left(x - \frac {\langle x , a ^ {*} \rangle}{\langle x _ {0} , a ^ {*} \rangle} x _ {0}\right), \frac {| \langle x , a ^ {*} \rangle |}{| \langle x _ {0} , a ^ {*} \rangle |} \alpha (x _ {0})\right).
\]

由此推出：存在 E 的一个基 \((a_{i})\) 以及一个实数的族 \((r_i)\)（均 \(>0\)），使得对所有 \(x = \sum_{i} \xi_{i} a_{i}\)，有 \(\alpha(x) = \sup_{i} (r_{i} |\xi_{i}|)\)。我们称 \(\alpha\) 关于基 \((a_{i})\) 具有标准形式。

\begin{enumerate}
\def\labelenumi{\alph{enumi})}
\setcounter{enumi}{1}
\tightlist
\item
  设 \(\alpha^{*}\) 是 E 的对偶 E* 上由 \(\alpha\) 通过下式典范地确定的范数
\end{enumerate}

\[
\alpha^ {*} (x ^ {*}) = \sup _ {x \neq 0} | \langle x, x ^ {*} \rangle | / \alpha (x);
\]

它是一个超范数。证明：对所有 \(x_0 \neq 0\)（在 \(\mathbf{E}\) 中），存在 \(x_0^* \in \mathbf{E}^*\)，使得 \(\alpha(x_0) = |\langle x_0, x_0^* \rangle| / \alpha^*(x_0^*)\)。

\begin{enumerate}
\def\labelenumi{\alph{enumi})}
\setcounter{enumi}{2}
\tightlist
\item
  设 \(\alpha, \beta\) 是 E 上任意两个超范数。证明：存在 E 的一个基，使 \(\alpha\) 与 \(\beta\) 关于该基都具有标准形式（考察一点 \(x_0 \in \mathrm{E} \setminus \{0\}\)，\(\alpha / \beta\) 在该点达到其最大值；然后利用 \(b\)），并对 \(\dim \mathrm{E}\) 施行归纳法）。
\end{enumerate}

\(d)\) 把 \(\mathfrak{N}_0(\mathrm{E})\)（\(\mathrm{E}\) 上全体超范数所成的集）视为 \(\mathfrak{N}(\mathrm{E})\)（习题 10）的子空间。证明 \(\mathfrak{N}_0(\mathrm{E})\) 在 \(\mathfrak{N}(\mathrm{E})\) 中是闭的。设 \(\alpha_0\) 是 \(\mathfrak{N}_0(\mathrm{E})\) 的一个元素；对每个 \(\alpha \in \mathfrak{N}_0(\mathrm{E})\) 以及 \(0 \leqslant t \leqslant 1\)，用 \(\mathrm{P}_{\alpha}(t)\) 表示诸 \(\beta \in \mathfrak{N}_0(\mathrm{E})\)（满足 \(\beta(x) \leqslant \alpha_0(x)^{1 - t}\alpha(x)^t\)，对所有 \(x \in \mathrm{E}\)）组成的集。证明 \(\mathrm{P}_{\alpha}(t)\) 非空，且 \(\pi_t^\alpha = \sup \mathrm{P}_{\alpha}(t)\) 是一个超范数。此外，映射 \((t, \alpha) \mapsto \pi_t^\alpha\)（\([0, 1] \times \mathfrak{N}_0(E)\) 到 \(\mathfrak{N}_0(E)\) 中的）是连续的，且满足 \(\pi_0^\alpha = \alpha_0\) 与 \(\pi_1^\alpha = \alpha\)（利用 \(c\)）。

* e) 设 A 是 K 的绝对值的环，m 是其极大理想，使得 \(k = A/m\) 是一个含 q 个元素的有限域 (CA, VI, § 5, No.~1, prop. 2)。对 E 上每个超范数 \(\alpha\)，记 \(X_{\alpha}\) 为 \(\log \alpha(x)\)（\(x \in E \setminus \{0\}\)）的值集在到商群 \(\mathbf{R}/(\mathbf{Z}, \log q)\) 中的典范映射下的像；它是一个至多含 \(n = \dim E\) 个元素的有限集（利用 a)）；这些元素的个数记作 \(r(\alpha)\)，称为 \(\alpha\) 的秩。证明 r 是 \(\mathfrak{N}_{0}(E)\) 到 N 中的一个下半连续映射，且集 \(\mathfrak{N}_{0}'(E)\)（由诸 \(\alpha\)（\(r(\alpha) = n\)）所组成）在 \(\mathfrak{N}_{0}(E)\) 中是开的且处处稠密（利用 a) 与 c)）。

\begin{enumerate}
\def\labelenumi{\alph{enumi})}
\setcounter{enumi}{5}
\tightlist
\item
  设 \(r(\alpha) = n\)；设 \((a_i)\) 是 E 的一个基，\(\alpha\) 关于它具有标准形式；证明：存在 \(\alpha\) 在 \(\mathfrak{N}_0'(\mathrm{E})\) 中的一个邻域 V，使得每个 \(\beta \in V\) 都具有关于 \((a_i)(\text{使用 } b)\) 的标准形式；由此推出，存在邻域 \(W \subset V\)（\(\alpha\) 的），它同胚于 \(R^n\) 中的一个开集。
\end{enumerate}

\(g)\) 对 E 的每个基 \((a_i)\)，证明：超范数 \(\alpha\)（具有关于 \((a_i)\) 的标准形式）所成的集在 \(\mathfrak{N}_0^{\prime}(\mathrm{E})\) 中是闭的。由此推出：若 \(\alpha \in \mathfrak{N}_0^{\prime}(\mathrm{E})\) 具有关于 \((a_i)\) 的标准形式，则含 \(\alpha\) 的连通分支（\(\mathfrak{N}_0^{\prime}(\mathrm{E})\) 中的）的所有元素亦然。

¶ 13) * 我们保留习题 12 的一般假设与记号。

\begin{enumerate}
\def\labelenumi{\alph{enumi})}
\item
  设 \(\mathbf{L}\) 是 \(\mathbf{E}\) 的一个维数为 \(n = \dim \mathbf{E}\) 的自由 A-子模。对所有 \(x \in \mathbf{E} \setminus \{0\}\)，诸 \(a \in \mathbf{A}\)（满足 \(ax \in \mathbf{L}\)）组成的集是 \(\mathbf{K}\) 的一个形如 \(\mathfrak{m}^h\) 的分式理想（\(h\) 为正整数或负整数）；置 \(\alpha(x) = q^h\) 及 \(\alpha(0) = 0\)，证明 \(\alpha\) 是 \(\mathbf{E}\) 上的一个超范数。称它为与自由 A-模 \(\mathbf{L}\) 相伴的超范数。
\item
  反之，若 \(\alpha\) 是 E 上的一个超范数，则集 \(L_{\alpha}\)（由诸 \(x \in E\)（满足 \(\alpha(x) \leqslant 1\)）组成）是一个维数为 \(n\) 的自由 A-模。若 \([\alpha]\) 是与 \(L_{\alpha}\) 相伴的范数，则有 \(\alpha \leqslant [\alpha] \leqslant q\alpha\)，且 \([\alpha]\) 是与自由 A-模相伴且 \(\geqslant \alpha\) 的诸范数的下确界。我们有 \([q\alpha] = q \cdot [\alpha]\)，且 \(\alpha(x) = \inf(q^{-t}[q^t\alpha](x))\)（对所有 \(x \in E\)），其中 \(t\) 取遍区间 {[}0, 1{]}。此外，函数 \(t \mapsto [q^t\alpha](x)\) 在该区间上是左连续的。
\item
  沿用同样的记号，证明：对 \(0 \leqslant t \leqslant 1\)，诸 \([q^{t}\alpha]\) 中至多有 n 个互不相同的超范数。反之，设 L 是与维数为 n 的自由 A-模相伴的超范数所成的集，\((\alpha_{t})_{0 \leqslant t \leqslant 1}\) 是 L 的一个递增的超范数族，且满足 \(\alpha_{1} = q\alpha_{0}\)。证明：存在 E 的一个基，使所有 \(\alpha_{t}\) 都关于它具有标准形式（若 \(u \in A\) 是赋值为 1 的一个元素，\(L_{t}\) 是诸 \(x \in E\)（满足 \(\alpha_{t}(x) \leqslant 1\)）组成的自由 A-模，则考察 k 上的向量空间 \(L_{t}/uL_{0}\)）。进一步推出：若对所有 \(x \in E\)，\(t \mapsto \alpha_{t}(x)\) 在 \([0, 1]\) 上左连续，则存在唯一的超范数 \(\alpha\)，使得 \(\alpha_{t} = [q^{t}\alpha]\)（对所有 \(t \in [0, 1]\)）。
\end{enumerate}

\(d\) ) 线性群 \(\mathbf{GL}(\mathrm{E})\) 连续地作用于 \(\mathfrak{N}_0(\mathrm{E})\)；证明它是真作用。对所有 \(\alpha \in \mathfrak{N}_0(\mathrm{E})\)，稳定子 \(S_{\alpha}\)（\(\alpha\) 在 \(\mathbf{GL}(\mathrm{E})\) 中的）是诸 \([q^t\alpha]\)（\(0\leqslant t\leqslant 1\)）的稳定子的交；由此推出 \(S_{\alpha}\) 是 \(\mathbf{GL}(\mathrm{E})\) 的一个开紧子群，从而每个 \(\alpha \in \mathfrak{N}_0(\mathrm{E})\) 的轨道都是 \(\mathfrak{N}_0(\mathrm{E})\) 的一个闭离散子空间。

\begin{enumerate}
\def\labelenumi{\alph{enumi})}
\setcounter{enumi}{4}
\tightlist
\item
  对每个超范数 \(\alpha \in \mathfrak{N}_0(\mathrm{E})\)，考察向量 \(k\)-空间 \(\mathrm{L}_t / u\mathrm{L}_0\) 的维数所成的递减序列，其中 \(\mathrm{L}_t\) 是诸 \(x \in \mathrm{E}\)（满足 \([q^t\alpha](x) \leqslant 1\)）组成的 A-模，而 \(t\) 从 0 变到 1；我们称该序列为 \(\alpha\) 的不变量序列。为使 \(\alpha\) 与 \(\beta\) 属于 \(\mathfrak{N}_0(\mathrm{E})\) 中同一轨道，必须且只需 \(\mathrm{X}_{\alpha} = \mathrm{X}_{\beta}\)（习题 12, e)）且 \(\alpha\) 与 \(\beta\) 的不变量序列相同（利用习题 12, b)）。f) 由 e) 推出：轨道空间 \(\mathfrak{N}_0(\mathrm{E}) / \mathbf{GL}(\mathrm{E})\) 同构于轨道空间 \(\mathbf{T}^n / \mathfrak{S}_n\)，其中对称群在 \(\mathbf{T}^n\) 上通过 \((z_1, ..., z_n) \mapsto (z_{\sigma(1)}, ..., z_{\sigma(n)})\) 右作用。*
\end{enumerate}

\begin{enumerate}
\def\labelenumi{\arabic{enumi})}
\setcounter{enumi}{13}
\item
  把第 2 小节与第 3 小节的结果推广到离散除环 K 上的这样的拓扑向量空间 E：E 中存在由均衡邻域组成的 0 的基本邻域系（即由满足 K.V = V 的邻域 V 组成的基本邻域系）。
\item
  设 \(\mathbf{E}\) 是一个维数 \(n\) 有限的赋范空间（在 \(\mathbf{R}\) 或 \(\mathbf{C}\) 上）。给对偶 \(\mathbf{E}^*\) 赋以由 \(\| x^* \| = \sup_{\| x \| \leqslant 1} |\langle x, x^* \rangle|\) 定义的范数 (GT, X, § 3.2)。证明：存在一个基 \((e_i)\)
\end{enumerate}

使 E 满足：若 \((e_{i}^{*})\) 是对偶基，则对所有 i 有 \(\|e_{i}\| = \|e_{i}^{*}\| = 1\)。（设 \((a_{i})\) 是由范数为 1 的向量组成的 E 的一个基；考察行列式 \(\det(\xi_{ij})\)（对每个由 n 个范数为 1 的向量 \(x_{i} = \sum_{j} \xi_{ij} a_{j}\) 组成的向量组），并考察使该行列式的绝对值达到最大的这样一个向量组。）

